\documentclass[12pt]{extarticle}
\title{Análisis de juegos}
\usepackage{gensymb}
\usepackage{graphicx}
\usepackage{amssymb}
\usepackage{amsthm}
\usepackage{amsmath}
\usepackage{xcolor}
\usepackage{tikz}
\usepackage{float}
\usetikzlibrary{positioning}
\usetikzlibrary{calc, shapes.geometric}
\usetikzlibrary{intersections}
\usepackage{subcaption}

\usepackage{xcolor}
\usepackage{tikz}
\usepackage{tkz-euclide}
\usetikzlibrary{arrows}
\usetikzlibrary{calc,through}
 \usetikzlibrary{math}
 
\tikzset{circle through 3 points/.style n args={3}{%
insert path={let    \p1=($(#1)!0.5!(#2)$),
                    \p2=($(#1)!0.5!(#3)$),
                    \p3=($(#1)!0.5!(#2)!1!-90:(#2)$),
                    \p4=($(#1)!0.5!(#3)!1!90:(#3)$),
                    \p5=(intersection of \p1--\p3 and \p2--\p4)
                    in },
at={(\p5)},
circle through= {(#1)}
}}
\usetikzlibrary{shapes.geometric}


\usepackage[shortlabels]{enumitem}
\setenumerate[0]{label=\alph*)}
\graphicspath{ {images/} }
\begin{document}

\newtheorem{ejer} {Reto}
\newtheorem{ejem} {Ejemplo}
\newtheorem{theorem} {Teorema}
\newtheorem*{demon} {Demostración}
\newtheorem*{pistas} {Pista}
\newtheorem*{pistas2} {Pista 2}
\newtheorem*{sol} {Solución}
\newtheorem{defin} {Definición}
\newenvironment{enumerar}{\begin{enumerate}\setlength{\parskip}{0pt}}
{\end{enumerate}}

\newcommand{\temas}[1]{}
\newcommand{\dificultad}[1]{}
\newcommand{\fuente}[1]{}
\newcommand{\curso}[1]{}
\newcommand{\comentarios}[1]{}



%%%%%%%%%%%%%%%%%%%%%%%%%%%%%%%%%%%%%%%%%%%%
% PROBLEMA 1
%%%%%%%%%%%%%%%%%%%%%%%%%%%%%%%%%%%%%%%%%%%%

\title{
% Título del problema
}

\temas{
% 1-4 tags apropiados, como 
% geometría, circunferencias, inversión
}

\dificultad{
% número del 1 al 6
%1: ejercicio técnico que sigue el algoritmo explicado en clase
%2: el 90% de alumnos listos lo resuelven en menos de 5 minutos
%3: el 60-70% de alumnos listos lo resuelven en menos de 10 minutos
%4: lo resuelve uno de cada 3-4 en 15-30 minutos
%5: lo resuelven menos del 10% de alumnos listos en más de media hora
%6: lo resuelven menos del 2% de alumnos listos
}

\fuente{
% indicar autor y procedencia, por ejemplo
% IMO 2026, Shapovalov V.
}

\curso{
% curso mínimo necesario para comprender y poder solucionar el problema sin explicaciones adicionales
% 1 Primaria, 2 Primaria,..., ! ESO,..., 1 BACH, 2 BACH
}

\comentarios{

}

\begin{ejer}
% el enunciado
\end{ejer}
\begin{pistas}
% las pistas
\end{pistas}

\begin{proof}
% la(s) solucion(es)
\end{proof}

\end{document}